\subsection{Same Cardinality, Finite, Countable, and Uncountable}

\begin{definition} \label{def:cardinality-equivalence} \hlt{Same Cardinality}\\
Two sets $A$ and $B$ have \hlt{same cardinality}, or are \hlt{equivalent}, if $\exists$ a bijection $f:A\to B$. Then $A\sim_{\#}B$.\\
The relation $\sim_{\#}$ is an equivalence relation: $A\sim_{\#}A$ by identity map, symmetry follows from Theorem \ref{thm:bijective-iff-invertible-function}, and transitivity follows by composing bijections.
\end{definition}

\begin{definition} \label{def:finite-countable-uncountable} \hlt{Finite, Countable, and Uncountable}\\
Set $J_0\defeq\emptyset$, and for $n\geq1$ set $J_n\defeq\{1,2,\ldots,n\}$. A set $A$ is
\begin{enumerate}[label=(\roman*)]
\setlength{\itemsep}{0pt}
\item \hlt{finite} if $A\sim_{\#} J_n$ for some $n\in\N$;
\item \hlt{infinite} if it is not finite;
\item \hlt{countable} if $A\sim_{\#} \N$;
\item \hlt{uncountable} if it is neither finite nor countable;
\item \hlt{at most countable} if it is finite or countable.
\end{enumerate}
Countable sets are also called enumerable or denumerable.
\end{definition}

\begin{proposition} \label{prop:finite-cardinality-well-defined} \hlt{Finite Cardinality is Well Defined}\\
If $A\sim_{\#}J_m$ and $A\sim_{\#}J_n$, then $m=n$. Thus $\#A=n$ is well defined for finite sets.
\begin{proof}
By transitivity of $\sim_{\#}$, it suffices to prove $J_m\sim_{\#}J_n\Rightarrow m=n$. Induct on $m$.\\
For $m=0$, only $J_0=\emptyset$ is bijective to $J_0$.\\
If $m,n\geq1$ and $f:J_m\to J_n$ is bijective, remove $m$ and its partner $f(m)$. The remaining sets are bijective and relabel as $J_{m-1}$ and $J_{n-1}$, so induction gives $m-1=n-1$, hence $m=n$.\\
\hlt{Structural Commentary}: Finite counting is invariant under bijection.\\
\hlt{Key Insight}: Remove one matched element and count down.
\end{proof}
\end{proposition}

\begin{definition} \label{def:sequence-as-enumeration} \hlt{Sequences as Enumerations}\\
A \hlt{sequence} in set $A$ is $x:\N\to A$, written $(x_n)$, where $x_n\defeq x(n)$. Terms of sequence need not be distinct.\\
If $A$ is countable, then its elements can be arranged as a sequence with distinct terms. Conversely, a sequence whose terms are distinct and whose range is $A$ gives a bijection $\N\to A$.
\end{definition}

\begin{example} \label{ex:integers-are-countable} \hlt{The Integers Are Countable}\\
The set $\Z$ is countable. One explicit enumeration is given by
\begin{align}
f(n)=
\begin{cases}
\frac{n}{2}, & n \text{ even},\\[2mm]
-\frac{n-1}{2}, & n \text{ odd},
\end{cases}
\qquad n\in\N. \nonumber
\end{align}
This lists $\Z$ as $0,1,-1,2,-2,3,-3,\ldots$.
\end{example}

\begin{theorem} \label{thm:infinite-subset-of-countable-is-countable} \hlt{Infinite Subsets of Countable Sets}\\
Every infinite subset of a countable set is countable.
\begin{proof}
Let $E\subseteq A$ be infinite, and enumerate $A$ by distinct terms $(x_n)$. Let $n_1<n_2<\cdots$ be exactly the indices for which $x_{n_k}\in E$, listed in increasing order. Then $k\mapsto x_{n_k}$ is a bijection from $\N$ onto $E$.\\
\hlt{Structural Commentary}: The proof keeps only the entries of the original enumeration that land in $E$.\\
\hlt{Key Insight}: An infinite subset of an enumeration inherits an enumeration by first appearances.
\end{proof}
\end{theorem}

\begin{proposition} \label{prop:subsets-and-images-at-most-countable} \hlt{Subsets and Images of At Most Countable Sets}\\
Let $A$ be at most countable.
\begin{enumerate}[label=(\roman*)]
\setlength{\itemsep}{0pt}
\item If $B\subseteq A$, then $B$ is at most countable.
\item If $f:A\to Y$ is a function, then $f(A)$ is at most countable.
\end{enumerate}
\begin{proof}
If $A$ is finite, both claims are immediate. Otherwise write $A=\{a_n:n\in\N\}$ with the $a_n$ distinct.
\begin{enumerate}[label=(\roman*)]
\setlength{\itemsep}{0pt}
\item If $B$ is infinite, let $I_B\defeq\{n:a_n\in B\}=\{n_1<n_2<\cdots\}$.\\
Then $k\mapsto a_{n_k}$ bijects $\N$ onto $B$; the finite case is trivial.
\item Let $I\defeq\{n:f(a_n)\notin\{f(a_m):m<n\}\}$. By (i), $I$ is at most countable, $n\mapsto f(a_n)$ bijects $I$ onto $f(A)$.
\end{enumerate}
\hlt{Structural Commentary}: Subsets select indices; images select first appearances.\\
\hlt{Key Insight}: Subsets and images preserve at-most-countability.
\end{proof}
\end{proposition}

\begin{figure}[H]
  \centering
  \begin{tikzpicture}[
  >=Stealth,
  every node/.style={font=\small},
  entry/.style={draw=gray!55, fill=gray!5, rounded corners=1pt, minimum width=0.62cm, minimum height=0.40cm},
  active/.style={draw=orange!60!black, fill=orange!10},
  scan/.style={->, blue!70!black, line width=0.7pt},
  guide/.style={orange!75!black, line width=0.65pt},
  output/.style={draw=gray!55, fill=gray!5, rounded corners=1pt, inner xsep=5pt, inner ysep=4pt}
]
  \draw[guide] (1.70,-0.65) -- (0.85,-1.30);
  \draw[guide] (2.55,-0.65) -- (0.85,-1.95);
  \draw[guide] (3.40,-0.65) -- (0.85,-2.60);

  \foreach \row/\name in {1/E_1,2/E_2,3/E_3,4/E_4} {
    \node[anchor=east] at (0,-0.65*\row) {$\name$};
    \foreach \col in {1,...,5} {
      \node[entry] (x-\row-\col) at (0.85*\col,-0.65*\row) {$x_{\row\col}$};
    }
    \node[anchor=west, gray!65] at (4.55,-0.65*\row) {$\cdots$};
  }

  \foreach \row/\col in {1/1,1/2,2/1,1/3,2/2,3/1,1/4,2/3,3/2,4/1} {
    \node[entry, active] at (0.85*\col,-0.65*\row) {$x_{\row\col}$};
  }

  \node[output, anchor=west] (list) at (5.25,-1.45)
    {$x_{11},\,x_{12},\,x_{21},\,x_{13},\,x_{22},\,x_{31},\,x_{14},\,\cdots$};
  \draw[scan] (4.55,-1.45) -- (list.west);
\end{tikzpicture}

  \caption{Diagonal scan of countable rows into one sequence, skipping repeats.}
  \label{fig:countable-union-diagonal-scan}
\end{figure}

\begin{theorem} \label{thm:countable-union-of-countable-sets}\label{thm:cardinality-main-results-source-todo} \hlt{Countable Union of Countable Sets}\\
Let $(E_n)$ be a sequence of countable sets and put $S\defeq\bigcup\limits_{n=1}^{\infty}E_n$. Then $S$ is countable.
\begin{proof}
Enumerate each $E_n$ as $(x_{nk})_{k\in\N}$.\\
Reading the array $(x_{nk})$ by diagonals, as in Figure \ref{fig:countable-union-diagonal-scan}, gives a sequence whose image is $S$.\\
Thus $S$ is at most countable by Proposition \ref{prop:subsets-and-images-at-most-countable}.\\
Since $E_1\subseteq S$ and $E_1$ is countable, $S$ is infinite; hence $S$ is countable.\\
\hlt{Structural Commentary}: The diagonal order turns a two-indexed list into a one-indexed list.\\
\hlt{Key Insight}: Countably many countable rows can still be scanned by one sequence.
\end{proof}
\end{theorem}

\begin{corollary} \label{cor:at-most-countable-union-at-most-countable} \hlt{At Most Countable Unions}\\
If $A$ is at most countable and $B_\alpha$ is at most countable for each $\alpha\in A$, then $\bigcup\limits_{\alpha\in A}B_\alpha$ is at most countable.
\begin{proof}
Let $T\defeq\bigcup\limits_{\alpha\in A}B_\alpha$. If $T=\emptyset$, we are done. Otherwise enumerate the nonempty sets among the $B_\alpha$ as $B_{\alpha_1},B_{\alpha_2},\ldots$, repeating indices if needed, and list each $B_{\alpha_j}$ as $(b_{jk})_{k\in\N}$.\\
A diagonal scan of $(b_{jk})$ has image $T$, so Proposition \ref{prop:subsets-and-images-at-most-countable} makes $T$ at most countable.\\
\hlt{Structural Commentary}: Finite pieces cause no difficulty; they can be embedded into countable rows.\\
\hlt{Key Insight}: At most countable indexing $+$ at most countable indexed sets gives at most countable total set.
\end{proof}
\end{corollary}

\begin{theorem} \label{thm:finite-tuples-over-countable-set-countable} \hlt{Finite Tuples over a Countable Set}\\
Let $A$ be countable. For each fixed $n\in\N$, the set of all $n$-tuples $(a_1,\ldots,a_n)$ with $a_k\in A$ is countable.
\begin{proof} Let $B_n\defeq \{(a_1,\ldots,a_n): a_k\in A\}$, $B_1=A$. If $B_{n-1}$ is countable, $B_n=\bigcup\limits_{b\in B_{n-1}}\{(b,a):a\in A\}$, each set $\{(b,a):a\in A\}$ is equivalent to $A$. Hence $B_n$ is countable by Theorem \ref{thm:countable-union-of-countable-sets}.\\
\hlt{Structural Commentary}: Build an $n$-tuple by appending one coordinate to an $(n-1)$-tuple.\\
\hlt{Key Insight}: Induction plus countable unions controls finite products.
\end{proof}
\end{theorem}

\begin{figure}[H]
  \centering
  \begin{tikzpicture}[
  x=0.82cm,
  y=0.82cm,
  >=Stealth,
  every node/.style={font=\small},
  point/.style={circle, fill=black, inner sep=1.25pt},
  active/.style={circle, draw=orange!60!black, fill=orange!12, line width=0.55pt, inner sep=2.2pt},
  guide/.style={orange!75!black, line width=0.65pt},
  scan/.style={->, blue!70!black, line width=0.7pt},
  output/.style={draw=gray!55, fill=gray!5, rounded corners=1pt, inner xsep=5pt, inner ysep=4pt}
]
  \foreach \i in {1,...,5} {
    \node[anchor=east] at (0,-0.42*\i) {$\i$};
    \node[anchor=south] at (0.47*\i,-0.1) {$\i$};
  }
  \node[anchor=south] at (2.65,-0.1) {$\cdots$};

  \draw[guide] (0.94,-0.42) -- (0.47,-0.84);
  \draw[guide] (1.41,-0.42) -- (0.47,-1.26);
  \draw[guide] (1.88,-0.42) -- (0.47,-1.68);

  \foreach \row in {1,...,5} {
    \foreach \col in {1,...,5} {
      \node[point] at (0.47*\col,-0.42*\row) {};
    }
  }

  \foreach \row/\col in {1/1,1/2,2/1,1/3,2/2,3/1,1/4,2/3,3/2,4/1} {
    \node[active] at (0.47*\col,-0.42*\row) {};
  }

  \node[output, anchor=west] (list) at (3.10,-1.05)
    {$(1,1),\,(1,2),\,(2,1),\,(1,3),\,(2,2),\,(3,1),\,\cdots$};
  \draw[scan] (2.70,-1.05) -- (list.west);
\end{tikzpicture}

  \caption{Diagonal enumeration of $\N\times\N$ by finite $i+j$ diagonals.}
  \label{fig:countable-diagonal-enumeration}
\end{figure}

\begin{corollary} \label{cor:countable-cartesian-products} \hlt{Countable Cartesian Products}\\
If $A,B$ are at most countable, then $A\times B$ is at most countable.\\
If $A,B$ are countable and nonempty, then $A\times B$ is countable. Hence $\N\times\N$ is countable.
\begin{proof}
For $a\in A$, $\{a\}\times B\sim B$, and $A\times B=\bigcup\limits_{a\in A}\bigl(\{a\}\times B\bigr)$, thus $A\times B$ at most countable (Corollary \ref{cor:at-most-countable-union-at-most-countable}). If $A,B$ are countable and $b_0\in B$, then $a\mapsto(a,b_0)$ injects $A$ into $A\times B$, so $A\times B$ is infinite; hence it is countable.\\
\hlt{Structural Commentary}: A product is a countable union of countable slices.\\
\hlt{Key Insight}: Two countable coordinates can still be scanned by one sequence.
\end{proof}
\end{corollary}

\begin{corollary} \label{cor:rationals-countable} \hlt{The Rationals Are Countable}\\
The set $\Q$ is countable.
\begin{proof}
Let $D\defeq\{(a,b)\in\Z\times\Z:a\neq0\}$ and $\pi:D\to\Q$ by $\pi(a,b)=b/a$.\\
Since $\Z$ is countable, Corollary \ref{cor:countable-cartesian-products} and Proposition \ref{prop:subsets-and-images-at-most-countable} make $D$ at most countable. Since $\pi(D)=\Q$, Proposition \ref{prop:subsets-and-images-at-most-countable} makes $\Q$ at most countable. As $\Z\subseteq\Q$, it is countable.\\
\hlt{Structural Commentary}: Repeated fractions only add repetitions to the list.\\
\hlt{Key Insight}: Countable pairs produce countably many rational values.
\end{proof}
\end{corollary}
